\chapter{Численные методы оптимизации}
\label{ch:opt}

\section{Метод касательных Ньютона и его приложения}
\label{sec:newton}

\subsection{От Ньютона до Канторовича}

Задачи отыскания корня уравнения $f(x)=0$ и поиска минимума функции
$g(x)$ относятся к разным «жанрам» школьной математики, однако
численно они решаются почти одним и тем же приёмом — итерационной
схемой, придуманной Исааком Ньютоном.

\begin{historybox}
В работе \emph{De analysi per aequationes numero terminorum infinitas}
(1669) и затем в \emph{Method of Fluxions} (рукопись 1671~г., издание
1736~г.) Ньютон предложил способ приближённого решения уравнений,
который позже Джозеф Рафсон в 1690~г. переписал в современной
рекуррентной форме. Поэтому в англоязычной литературе метод часто
называют \emph{Newton–Raphson}.

Существенный вклад в строгое исследование метода внёс советский
математик \textbf{Леонид Витальевич Канторович} (1912–1986),
лауреат Нобелевской премии по экономике 1975~г. (за теорию
оптимального распределения ресурсов). В 1948~г. он опубликовал
теорему, носящую теперь его имя: она впервые дала универсальные
\emph{достаточные условия сходимости} метода Ньютона, причём не
только для уравнений с одной переменной, но и для уравнений в
бесконечномерных банаховых пространствах.
\end{historybox}

\paragraph{Геометрическая идея.} Пусть требуется найти корень
уравнения~$f(x)=0$, и нам известна точка $x_k$, близкая к корню.
Заменим график $y=f(x)$ его \emph{касательной} в точке $(x_k, f(x_k))$.
Касательная — это прямая, и её пересечение с осью $x$ найти просто.
Точку этого пересечения и объявим следующим приближением~$x_{k+1}$.
Затем повторим. Так получается метод касательных~(рис.~\ref{fig:newton-geom}).

\begin{figure}[ht]
\centering
\includegraphics[width=0.85\linewidth]{newton_geometry.pdf}
\caption{Три итерации метода Ньютона для уравнения
$x^{3}-2x-5=0$. Каждая касательная (пунктир) пересекает ось абсцисс
в точке, которая берётся в качестве следующего приближения
(квадрат). После трёх шагов корень найден с точностью лучше $10^{-6}$.}
\label{fig:newton-geom}
\end{figure}

Уравнение касательной к графику $y=f(x)$ в точке $x_k$ имеет вид
$y=f(x_k)+f'(x_k)\,(x-x_k)$.
Полагая $y=0$ и выражая $x$, получаем
\begin{equation}\label{eq:newton-root}
\boxed{\;x_{k+1}=x_k-\dfrac{f(x_k)}{f'(x_k)}\;}\qquad k=0,1,2,\dots
\end{equation}

Это и есть \emph{итерационная формула Ньютона}. Точка $x_0$ выбирается
исследователем; от её выбора может зависеть, сойдётся ли метод
вообще.

\subsection{Метод Ньютона для задач оптимизации}

Пусть теперь нужно найти минимум функции $g$. В любой точке минимума
$x^{\star}$ выполнено $g'(x^{\star})=0$. Значит, минимизация
дифференцируемой функции — это решение уравнения
\[
g'(x)=0.
\]
Применим к нему формулу~\eqref{eq:newton-root} с $f=g'$ (и тогда $f'=g''$):
\begin{equation}\label{eq:newton-opt}
\boxed{\;x_{k+1}=x_k-\dfrac{g'(x_k)}{g''(x_k)}\;}
\end{equation}

Эту же формулу можно получить иначе — и это ценная для всей теории
оптимизации точка зрения. Разложим $g$ в окрестности~$x_k$ в
\textbf{ряд Тейлора}:
\begin{equation}\label{eq:taylor}
g(x)\;=\;g(x_k)\;+\;g'(x_k)(x-x_k)\;+\;\tfrac12 g''(x_k)(x-x_k)^2
   \;+\;o\!\bigl((x-x_k)^2\bigr).
\end{equation}

Откуда берётся такое разложение? Идея проста: рядом с точкой $x_k$
функция $g$ ведёт себя «почти как многочлен». Линейное приближение
$g(x_k)+g'(x_k)(x-x_k)$ — это уже знакомая нам касательная к
графику. Если мы хотим добиться большей точности, надо учесть и
\emph{кривизну} графика — а она как раз кодируется второй производной
$g''$. Так возникает квадратичный многочлен в~\eqref{eq:taylor}.

Отбросим бесконечно малое слагаемое и обозначим оставшийся квадратичный
многочлен через
\[
q_k(x)\;=\;g(x_k)+g'(x_k)(x-x_k)+\tfrac12 g''(x_k)(x-x_k)^2.
\]
Минимизировать $q_k$ умеет любой школьник: это квадратное уравнение,
решаемое в одну строку. Условие $q_k'(x)=0$ даёт $g'(x_k)+g''(x_k)(x-x_k)=0$,
откуда $x=x_k-g'(x_k)/g''(x_k)$, что в точности совпадает
с~формулой~\eqref{eq:newton-opt}.

\begin{notebox}[title={«Разделяй и властвуй»}]
Метод Ньютона для оптимизации — классический пример принципа
\emph{разделяй и властвуй}: вместо того чтобы решать сложную задачу
минимизации произвольной функции $g$ целиком, мы заменяем её
на каждом шаге простой подзадачей — минимизацией параболы $q_k$.
Парабола подбирается так, чтобы локально (в окрестности $x_k$)
имитировать поведение $g$, и её минимум вычисляется явно.
\end{notebox}

\begin{figure}[ht]
\centering
\includegraphics[width=0.78\linewidth]{taylor_parabola.pdf}
\caption{Метод Ньютона как минимизация локальной квадратичной модели
$q_k$ (пунктир): следующая точка $x_{k+1}$ — это вершина параболы,
касающейся графика $g$ в точке $x_k$.}
\label{fig:taylor}
\end{figure}

\subsection{Пример 1. Извлечение квадратного корня}
\label{ssec:heron}

Пусть $a>0$. Вычислим $\sqrt{a}$, применяя метод Ньютона
к уравнению
\[
f(x)=x^{2}-a=0.
\]
Тогда $f'(x)=2x$, и формула~\eqref{eq:newton-root} принимает вид
\begin{equation}\label{eq:heron}
\boxed{\;x_{k+1}=\dfrac12\!\left(x_k+\dfrac{a}{x_k}\right)\;}
\end{equation}

Это — знаменитый \textbf{метод Герона}, или вавилонский метод
квадратного корня.

\begin{historybox}
Формулу~\eqref{eq:heron} в явном виде описал \emph{Герон Александрийский}
(I век н.~э.) в трактате «Метрика». Однако эту итерацию умели
применять задолго до него: на вавилонской глиняной табличке
\textsf{YBC~7289} (примерно XVIII–XVI вв. до н.~э.) сохранилось
вычисление $\sqrt{2}$ с точностью до шестого десятичного знака —
именно с помощью среднего арифметического и среднего гармонического,
что эквивалентно~\eqref{eq:heron}. Получается, что одна из
итерационных схем «пережила» три с половиной тысячи лет и
по сей день используется в библиотеках вычисления корня.
\end{historybox}

\subsubsection*{Сходимость через сжимающее отображение}

Запишем~\eqref{eq:heron} как $x_{k+1}=T(x_k)$, где
$T(x)=\tfrac12(x+a/x)$. Точка $\sqrt{a}$ — \emph{неподвижная} для $T$:
$T(\sqrt{a})=\sqrt{a}$.

\begin{theorem}{Сходимость метода Герона}{heron}
Для любого $x_0>0$ последовательность~\eqref{eq:heron} монотонно
убывает (начиная с $k=1$) и сходится к $\sqrt{a}$. Более того, для
ошибки $e_k=x_k-\sqrt{a}$ выполнено
\begin{equation}\label{eq:heron-err}
e_{k+1}\;=\;\dfrac{e_k^2}{2x_k}\,, \qquad
\text{и при }k\geq 1\colon\quad
0\le e_{k+1}\le \dfrac{e_k^{2}}{2\sqrt{a}}.
\end{equation}
\end{theorem}

\begin{proof}
\textbf{Шаг 1 (нижняя оценка).} По неравенству о средних
(AM--GM) для $x_0>0$:
$x_1=\tfrac12(x_0+a/x_0)\geq \sqrt{x_0\cdot a/x_0}=\sqrt{a}$.
Значит, начиная с $x_1$, все члены лежат в~$[\sqrt{a},\infty)$.

\textbf{Шаг 2 (рекуррентность для ошибки).}
\[
e_{k+1}=x_{k+1}-\sqrt{a}=\dfrac{x_k+a/x_k}{2}-\sqrt{a}
=\dfrac{x_k^{2}-2\sqrt{a}\,x_k+a}{2x_k}=\dfrac{(x_k-\sqrt{a})^2}{2x_k}=\dfrac{e_k^{2}}{2x_k}.
\]
При $x_k\geq\sqrt{a}$ имеем $\tfrac{1}{2x_k}\leq\tfrac{1}{2\sqrt{a}}$,
откуда вторая часть~\eqref{eq:heron-err}.

\textbf{Шаг 3 (сжатие).} Производная $T'(x)=\tfrac12(1-a/x^{2})$.
На луче $x\geq\sqrt{a}$ выполнено $T'(x)\in[0,\tfrac12)$. Значит,
$T$~— \emph{сжимающее отображение} (с константой~$\tfrac12$) на
$[\sqrt{a},\infty)$. По теореме Банаха о неподвижной точке
последовательность $x_k=T^{k}(x_1)$ сходится к единственной
неподвижной точке~$\sqrt{a}$.
\end{proof}

Из~\eqref{eq:heron-err} видно, что после первой итерации ошибка
\emph{квадрируется}: каждые два десятичных знака удваиваются за
один шаг. Поэтому при разумном $x_0$ хватает 4--5 итераций, чтобы
получить машинную точность.

\begin{example}{Расстояние до горизонта}{horizon}
Стоящий человек ростом $h=1{,}70$~м смотрит на море. До какой точки
далеко он видит? Из теоремы Пифагора (рис.~\ref{fig:horizon}) для прямой,
касательной к поверхности Земли:
\[
d^{2}+R^{2}=(R+h)^{2},\qquad
d=\sqrt{2Rh+h^{2}}\;\approx\;\sqrt{2Rh}.
\]
При $R=6371$~км, $h=1{,}7\cdot 10^{-3}$~км:
\[
2Rh\;\approx\;2\cdot 6371\cdot 0{,}0017\;=\;21{,}66\;\text{(км$^2$)}.
\]
Применим к этому числу метод Герона со стартом $x_0=5$:
\[
\begin{array}{l|l}
k & x_k \\\hline
0 & 5{,}000\,000 \\
1 & \tfrac{1}{2}(5+21{,}66/5)=4{,}666\,000 \\
2 & 4{,}654\,463 \\
3 & 4{,}654\,460
\end{array}
\]
Получаем $d\approx 4{,}65$~км — расстояние, которое запоминается
любому, кто хоть раз стоял у моря.
\end{example}

\begin{figure}[ht]
\centering
\begin{tikzpicture}[scale=0.9, line cap=round]
  % Земля
  \draw[thick, accent] (0,0) circle (2.6);
  \node[accent!80!black] at (0,-0.4) {Земля};
  % центр
  \fill (0,0) circle (1.2pt) node[below right=1pt] {$O$};
  % точка наблюдателя
  \coordinate (A) at (0,2.6);
  \coordinate (P) at (0,3.05);          % человек на высоте h
  \draw[very thick, accent2] (A) -- (P);
  \fill (P) circle (1.2pt) node[left] {$P$};
  \fill (A) circle (1pt) node[left=1pt] {$A$};
  % касательная: из P касается окружности в точке T
  % используем геометрию: |OP|=R+h, расстояние от O до T равно R, угол OTC=90.
  \pgfmathsetmacro{\Rh}{3.05}
  \pgfmathsetmacro{\R}{2.6}
  \pgfmathsetmacro{\dlen}{sqrt(\Rh*\Rh-\R*\R)}    % длина касательной
  \pgfmathsetmacro{\sina}{\dlen/\Rh}
  \pgfmathsetmacro{\cosa}{\R/\Rh}
  \coordinate (T) at ({\R*\sina}, {\R*\cosa});
  \draw[thick, accent2] (P) -- (T);
  \draw[dashed] (0,0) -- (T);
  \fill (T) circle (1pt) node[right=2pt] {$T$};
  % прямой угол в T
  \draw ($(T)+({-0.18*sin(40)},{-0.18*cos(40)})$) -- 
        ++({0.18*cos(40)},{-0.18*sin(40)}) --
        ++({0.18*sin(40)},{0.18*cos(40)});
  % метки
  \node[accent2] at (0.45,2.8) {$h$};
  \node at (0.85,1.45) {$d$};
  \node[rotate=15] at (-0.7,1.4) {$R$};
\end{tikzpicture}
\caption{Расстояние до горизонта $d=\sqrt{(R+h)^2-R^2}\approx\sqrt{2Rh}$.}
\label{fig:horizon}
\end{figure}

\subsection{Пример 2. Вычисление обратного числа}

Что, если на нашем процессоре нет аппаратного деления, и нужно
вычислить $1/a$ при $a>0$ (только сложениями и умножениями)?
Применим метод Ньютона к уравнению
\[
f(x)=\dfrac{1}{x}-a=0.
\]
Тогда $f'(x)=-1/x^{2}$, и подстановка в~\eqref{eq:newton-root} даёт
\[
x_{k+1}=x_k-\dfrac{1/x_k-a}{-1/x_k^{2}}
       =x_k+x_k^{2}\bigl(1/x_k-a\bigr)
       =x_k\bigl(2-a x_k\bigr).
\]
То есть
\begin{equation}\label{eq:ns}
\boxed{\;x_{k+1}=x_k\,(2-a x_k)\;}
\end{equation}

Эта схема замечательна тем, что в ней нет ни одного деления — только
сложение и умножение. Поэтому именно она исторически использовалась
в первых ЭВМ для реализации операции деления на аппаратном уровне.

\subsubsection*{Анализ сходимости}

Введём «относительную невязку» $u_k=1-a x_k$. Тогда
\[
u_{k+1}=1-a x_{k+1}=1-a x_k(2-ax_k)=1-2 a x_k+(a x_k)^{2}=(1-a x_k)^{2}=u_k^{2}.
\]
Таким образом
\begin{equation}\label{eq:ns-error}
u_k=u_0^{\,2^{k}}\quad\Longrightarrow\quad u_k\to 0\;\Leftrightarrow\;|u_0|<1.
\end{equation}

Условие $|1-a x_0|<1$ эквивалентно $0<x_0<2/a$. \emph{Областью
сходимости} метода является интервал $\bigl(0,\,2/a\bigr)$, и
он \textbf{конечен}: вне этого интервала ($x_0>2/a$ или $x_0<0$)
итерации расходятся. Это принципиальное отличие от метода Герона,
который сходится для любого $x_0>0$.

\begin{notebox}[title={Метод Ньютона--Шульца и оптимизатор Muon}]
Если в~\eqref{eq:ns} заменить число $x_k$ на квадратную матрицу $X_k$,
а число $a$ — на матрицу $A$, получим знаменитую матричную итерацию
\textbf{Ньютона--Шульца}: $X_{k+1}=X_k(2I-A X_k)$. Она вычисляет
$A^{-1}$ без обращения матриц на каждом шаге — лишь умножениями.

Аналогичные полиномиальные итерации применяются для приближённой
\emph{ортогонализации} матриц. Именно так устроен оптимизатор
\textbf{Muon}, появившийся в 2024–2025~гг.\ и используемый при
обучении современных больших языковых моделей: в нём матрица
градиента на каждом шаге быстро ортогонализуется несколькими
итерациями Ньютона--Шульца — это намного дешевле, чем точное
сингулярное разложение.
\end{notebox}

\subsection{Реализация на Python и численный эксперимент}

Запишем универсальную функцию, реализующую метод Ньютона
по формуле~\eqref{eq:newton-root}, и применим её к двум разобранным
примерам. Комментарии в коде даны по-английски (как принято в
большинстве библиотек): \texttt{f}, \texttt{df} — сама функция и
её производная; \texttt{x0} — начальное приближение; \texttt{n\_iter} —
максимальное число итераций; \texttt{tol} — точность остановки;
функция возвращает массив \texttt{xs} с историей итераций.

\begin{lstlisting}[label={lst:newton}]
import numpy as np
import matplotlib.pyplot as plt


def newton(f, df, x0, n_iter=20, tol=1e-15):
    """Newton's tangent method.

    Parameters
    ----------
    f, df  : callables f(x) and f'(x).
    x0     : initial guess.
    n_iter : maximum number of iterations.
    tol    : stopping tolerance on |f(x)|.

    Returns
    -------
    Array of iterates, including x0.
    """
    xs = [x0]
    x = x0
    for _ in range(n_iter):
        fx = f(x)
        if abs(fx) < tol:
            break
        x = x - fx / df(x)
        xs.append(x)
    return np.array(xs)


# ---- Example 1. Square root of a ---------------------------------
a = 2.0
xs_sqrt = newton(lambda x: x**2 - a,
                 lambda x: 2 * x,
                 x0=50.0)          # very poor initial guess
err_sqrt = np.abs(xs_sqrt - np.sqrt(a))

# ---- Example 2. Reciprocal 1/a -----------------------------------
a = 3.0
xs_inv  = newton(lambda x: 1.0 / x - a,
                 lambda x: -1.0 / x**2,
                 x0=0.20)          # inside the convergence radius (0, 2/a)
err_inv = np.abs(xs_inv - 1.0 / a)

# ---- Convergence plot in log scale -------------------------------
plt.semilogy(err_sqrt, 'o-', label=r'sqrt: $|x_k-\sqrt{2}|$')
plt.semilogy(err_inv,  's-', label=r'1/a:  $|x_k-1/3|$')
plt.xlabel('iteration $k$'); plt.ylabel('error')
plt.legend(); plt.grid(True); plt.show()
\end{lstlisting}

\noindent\small\textit{Листинг~2.1. Универсальный метод касательных
Ньютона и его применение.}\normalsize

\paragraph{Что показывает эксперимент.}

На рис.~\ref{fig:sqrt-conv} приведены кривые сходимости метода Герона
для $\sqrt{2}$ при разных стартах $x_0\in\{0{,}1;\,1;\,5;\,50\}$.
Все четыре кривые сходятся к нулю, причём поведение типичное для
квадратичной сходимости: после нескольких «согласующих» шагов
число верных знаков удваивается за итерацию. Это иллюстрирует
\textbf{глобальную сходимость} метода Герона: годится \emph{любое}
начальное приближение из $(0,\infty)$.

\begin{figure}[ht]
\centering
\includegraphics[width=0.85\linewidth]{sqrt_convergence.pdf}
\caption{Глобальная сходимость метода Герона для $\sqrt{2}$ из четырёх
очень разных стартовых точек. Шкала ошибки логарифмическая.}
\label{fig:sqrt-conv}
\end{figure}

Совсем другое поведение демонстрирует метод Ньютона--Шульца для
$1/a$ при $a=3$ (рис.~\ref{fig:rec-conv}): область сходимости — это
интервал $(0,\,2/3)\approx(0,\,0{,}667)$.
Стартуя \emph{внутри} этого интервала, метод сходится квадратично
(нижние три кривые); вблизи правой границы (например, $x_0=0{,}66$) — сходится,
но «впритык»; а уже при $x_0=0{,}70$ ошибка взрывообразно растёт:
$|u_k|=|u_0|^{2^{k}}\to\infty$ в полном согласии с
\eqref{eq:ns-error}. Это \textbf{локальная сходимость}.

\begin{figure}[ht]
\centering
\includegraphics[width=0.85\linewidth]{reciprocal_convergence.pdf}
\caption{Локальная сходимость метода Ньютона--Шульца для $1/3$.
Область сходимости — интервал $(0,\,2/3)$; вне его метод расходится.}
\label{fig:rec-conv}
\end{figure}

\subsection{Достаточные условия сверхлинейной сходимости\texorpdfstring{$^\star$}{*}}
\label{ssec:thm-superlinear}

Сравнение двух примеров наводит на вопрос: \emph{когда именно}
метод Ньютона сходится так быстро, как мы наблюдаем? Ответ
даёт следующая теорема, восходящая к Канторовичу.

\begin{definition}{Сверхлинейная и квадратичная сходимости}{rates}
Говорят, что последовательность $\{x_k\}$ сходится к $x^{\star}$
со \textbf{сверхлинейной скоростью}, если
\[
\lim_{k\to\infty}\dfrac{|x_{k+1}-x^{\star}|}{|x_k-x^{\star}|}=0.
\]
Если же выполнено более сильное условие
$|x_{k+1}-x^{\star}|\leq M |x_k-x^{\star}|^{2}$ для некоторой
константы $M>0$, то сходимость называют \textbf{квадратичной}.
\end{definition}

\begin{definition}{Сильно выпуклая функция}{strconvex}
Функция $g\colon\R\to\R$ называется \textbf{сильно выпуклой} с
константой $\mu>0$, если $g\in C^{2}$ и $g''(x)\geq\mu$ для всех $x$.
Сильная выпуклость гарантирует существование и единственность
точки минимума и, что важно для нас, отделённость второй производной
от нуля.
\end{definition}

\begin{theorem}{Сверхлинейная сходимость метода Ньютона}{newton-rate}
Пусть выполнено хотя бы одно из условий:
\begin{enumerate}\itemsep0pt
  \item[(а)] (\emph{уравнение}) $f\in C^{2}$ в окрестности точки $x^{\star}$,
        $f(x^{\star})=0$ и $f'(x^{\star})\neq 0$;
  \item[(б)] (\emph{сильно выпуклая оптимизация}) $g\in C^{3}$,
        $g$~— сильно выпуклая на интервале, содержащем точку
        минимума $x^{\star}$, и применяется итерация~\eqref{eq:newton-opt}.
\end{enumerate}
Тогда существуют $\delta>0$ и $M>0$ такие, что для всякого
$x_0\in(x^{\star}-\delta,\,x^{\star}+\delta)$ итерации Ньютона
определены и
\begin{equation}\label{eq:rate}
|x_{k+1}-x^{\star}|\;\leq\;M\,|x_k-x^{\star}|^{2}.
\end{equation}
В частности, сходимость квадратичная (а значит, сверхлинейная).
\end{theorem}

\begin{proof}
Достаточно доказать случай (а): задача (б) сводится к нему,
если положить $f=g'$, тогда условие $f'(x^{\star})\neq 0$
эквивалентно $g''(x^{\star})\geq\mu>0$.

Разложим $f$ по формуле Тейлора с остаточным членом в форме
Лагранжа в окрестности~$x_k$ и подставим $x=x^{\star}$:
\[
0\;=\;f(x^{\star})\;=\;f(x_k)+f'(x_k)(x^{\star}-x_k)+\tfrac{f''(\xi_k)}{2}(x^{\star}-x_k)^{2},
\]
где $\xi_k$~— некоторая точка между $x_k$ и $x^{\star}$. Разделим
обе части на $f'(x_k)$ (по непрерывности эта величина не равна
нулю в малой окрестности $x^{\star}$):
\[
0\;=\;\dfrac{f(x_k)}{f'(x_k)}+(x^{\star}-x_k)+\dfrac{f''(\xi_k)}{2 f'(x_k)}(x^{\star}-x_k)^{2}.
\]
Из определения $x_{k+1}=x_k-f(x_k)/f'(x_k)$ получаем
$\frac{f(x_k)}{f'(x_k)}=x_k-x_{k+1}$. Подставляя:
\[
0=(x_k-x_{k+1})+(x^{\star}-x_k)+\dfrac{f''(\xi_k)}{2 f'(x_k)}(x^{\star}-x_k)^{2},
\]
откуда после очевидных переобозначений
\begin{equation}\label{eq:newton-recur}
x_{k+1}-x^{\star}\;=\;\dfrac{f''(\xi_k)}{2 f'(x_k)}\,(x_k-x^{\star})^{2}.
\end{equation}
Это и есть искомая \emph{рекуррентность для расстояния}.

Теперь подберём константы. По непрерывности $f'$ и $f''$ существует
$\delta_1>0$ и числа $L$, $m>0$ такие, что
$|f''(\xi)|\leq L$ и $|f'(x)|\geq m$ для всех $\xi,x\in
[x^{\star}-\delta_1,\,x^{\star}+\delta_1]$.
Положим $M=L/(2m)$ и $\delta=\min\!\bigl(\delta_1,\,\tfrac{1}{2M}\bigr)$.
Тогда из~\eqref{eq:newton-recur} при $|x_k-x^{\star}|<\delta$:
\[
|x_{k+1}-x^{\star}|\leq M|x_k-x^{\star}|^{2}\leq M\delta\cdot|x_k-x^{\star}|\leq\tfrac12|x_k-x^{\star}|<\delta.
\]
Значит, последовательность не покидает $\delta$-окрестность и
оценка~\eqref{eq:rate} выполнена для всех $k$.
\end{proof}

\paragraph{Где условия выполняются для наших примеров?}

\textbf{Метод Герона.} Здесь $f(x)=x^{2}-a$, $f'(x)=2x$,
$f''(x)=2$. В точке $x^{\star}=\sqrt{a}>0$ выполнено $f'(x^{\star})=2\sqrt{a}\neq 0$.
\emph{Условия теоремы выполняются на всём луче $(0,\infty)$,
не содержащем ноль}, и константа в~\eqref{eq:rate} принимает вид
$M=\frac{1}{2\sqrt{a}}$ — в полном согласии с
явным выражением~\eqref{eq:heron-err}, полученным элементарно.

\textbf{Метод Ньютона--Шульца.} Здесь $f(x)=1/x-a$, $f'(x)=-1/x^{2}$,
$f''(x)=2/x^{3}$. В точке $x^{\star}=1/a$ имеем $f'(x^{\star})=-a^{2}\neq 0$,
и теорема применима в малой окрестности этой точки. Однако, в
отличие от метода Герона, оценка снизу $|f'(x)|\geq m>0$ \emph{ломается}
при $x\to\infty$: производная стремится к нулю. Именно поэтому
радиус сходимости в этом примере конечен и равен $1/a$ слева и справа
от $x^{\star}$, что в сумме даёт интервал $(0,\,2/a)$, найденный нами
ранее аналитически.

\medskip

Таким образом, два простейших примера полностью укладываются в
доказанную теорему: метод Ньютона сходится сверхлинейно \emph{там,
где} производная регулярна и отделена от нуля. Глобальную же
сходимость, как у Герона, приходится доказывать дополнительно — например,
ссылаясь на структуру функции (выпуклость) и теорему Банаха о
сжимающем отображении, как мы и сделали в разделе~\ref{ssec:heron}.

\begin{tasksbox}
\begin{enumerate}\itemsep1pt
  \item Применяя метод Ньютона к уравнению $f(x)=x^{n}-a=0$, получите
        итерационную формулу для $\sqrt[n]{a}$. Покажите, что для $n=3$
        она имеет вид $x_{k+1}=\frac{2 x_k+a/x_k^{2}}{3}$.
  \item Докажите, что для метода Герона из любого $x_0>0$ погрешность
        $e_1$ всегда положительна, после чего последовательность
        монотонно убывает. (Указание: используйте AM--GM.)
  \item Для метода Ньютона--Шульца при $a=2$ найдите наибольшее $x_0$,
        при котором итерации ещё сходятся. Сравните численно скорость
        сходимости при $x_0=0{,}1$ и $x_0=0{,}9$.
  \item$^{\star}$ Сформулируйте и докажите аналог
        теоремы~\ref{thm:newton-rate} для метода Ньютона минимизации
        многомерной сильно выпуклой функции $g\colon\R^{n}\to\R$.
        Какую роль играет матрица Гессе $\nabla^{2}g$?
  \item$^{\star}$ Прочитайте о теореме Канторовича в учебнике
        \href{https://en.wikipedia.org/wiki/Kantorovich_theorem}{«Kantorovich theorem»}
        и сформулируйте её в одну фразу: какие именно три величины
        должны быть «в нужном соотношении» для гарантии сходимости?
\end{enumerate}
\end{tasksbox}
